\begin{figure}[!b]
\centering
\begin{tikzpicture}[>=Stealth, every node/.style={font=\small},
 component/.style={draw,rounded corners,minimum width=2.25cm,minimum height=.85cm,align=center}]
\node[component,fill=black!4] (B) at (0,0) {$B$ near $0$\\mass $1/3$};
\node[component,draw=blue!60!black,fill=blue!4] (Y1) at (3.8,0)
 {$Y_1(B)$ near $e_1$\\mass $1/6$};
\node[component,draw=blue!60!black,fill=blue!4] (W1) at (-3.8,0)
 {$W_1(B)$ near $-e_1$\\mass $1/6$};
\node[component,draw=orange!60!black,fill=orange!5] (Y2) at (0,1.9)
 {$Y_2(B)$ near $e_2$\\mass $1/6$};
\node[component,draw=orange!60!black,fill=orange!5] (W2) at (0,-1.9)
 {$W_2(B)$ near $-e_2$\\mass $1/6$};
\draw[->,blue!60!black] (B) -- node[above] {$Y_1$} (Y1);
\draw[->,blue!60!black] (B) -- node[above] {$W_1$} (W1);
\draw[->,orange!60!black] (B) -- node[right] {$Y_2$} (Y2);
\draw[->,orange!60!black] (B) -- node[right] {$W_2$} (W2);
\end{tikzpicture}
\caption{Schematic positions of the five disjoint components, with their
marginal masses. Every optimal plan is concentrated on triples joining
a central point to its two images with the same subscript. The arrows
denote local diffeomorphisms, not straight transport trajectories.
For each indicated map \(H\) and Borel set \(S\subset B\),
\(\mu(H(S))=\mu(S)/2\).}
\label{fig:components}
\end{figure}
